\subsection{Error Handling Mechanism}
\label{subsec:error-handling}

\subsubsection{Overview}
The error handling mechanism provides consistent and maintainable error management across all layers of the application. It is built around a registry pattern that decouples error detection and response logic from 
the middleware layer, allowing new error types to be added without modifying any existing code.

\subsubsection{Core Abstractions}

\paragraph{The IErrorHandler Interface}
The foundation of the system is the \texttt{IErrorHandler} interface, which defines a unified contract that all error handlers must fulfill:

\begin{itemize}
    \item \texttt{canHandle(err: unknown): boolean} — determines whether this handler is responsible for the given error.
    \item \texttt{handle(err: unknown, res: Response): void} — sends the appropriate HTTP response for the error.
\end{itemize}

The use of \texttt{unknown} as the type for \texttt{err} in both methods ensures that every concrete handler fully honors the interface contract. Type narrowing happens internally inside each handler's \texttt{handle()} method, where the cast is safe because \texttt{canHandle()} has already confirmed the error type.

\paragraph{The AppError Class}
\texttt{AppError} is a custom error class that extends the native JavaScript \texttt{Error}. It is used throughout the application to represent known, expected failures such as resource not found or unauthorized access. It carries three properties:

\begin{itemize}
    \item \texttt{statusCode} — the HTTP status code for the response.
    \item \texttt{status} — derived automatically as \texttt{fail} for 4xx codes and \texttt{error} for 5xx codes.
    \item \texttt{payload} — optional additional data attached to the response.
\end{itemize}

\subsubsection{The Handler Registry Pattern}

\paragraph{Structure}
Each error category is encapsulated in its own self-contained handler class. The system includes the following handlers:

\begin{itemize}
    \item \texttt{PrismaKnownRequestErrorHandler} — handles known Prisma errors identified by error codes such as unique constraint violations (P2002), record not found (P2025), and connection errors (P1001, 
    P1002).
    \item \texttt{PrismaValidationErrorHandler} — handles malformed queries sent to Prisma.
    \item \texttt{PrismaUnknownRequestErrorHandler} — handles unexpected Prisma request failures.
    \item \texttt{PrismaRustPanicErrorHandler} — handles critical Prisma engine panics.
    \item \texttt{PrismaInitializationErrorHandler} — handles database initialization failures, with specific internal logging per failure cause.
    \item \texttt{JwtErrorHandler} — handles JWT errors including invalid signatures, expired tokens, and tokens used before their valid time.
    \item \texttt{BcryptErrorHandler} — handles errors arising from invalid arguments passed to hashing and comparison operations.
    \item \texttt{NodemailerErrorHandler} — handles SMTP connection and authentication errors with specific internal logging per error code.
    \item \texttt{AppErrorHandler} — handles all intentionally thrown \texttt{AppError} instances from application code.
    \item \texttt{UnexpectedErrorHandler} — a catch-all fallback whose \texttt{canHandle()} always returns \texttt{true}. It must always be the last entry in the registry.
    \item \texttt{CloudinaryErrorHandler} — handles Cloudinary storage service errors such as invalid file uploads, authentication failures, missing resources, and API rate limiting.
    \item \texttt{LLMErrorHandler} — handles large language model (LLM) API errors such as quota exhaustion, rate limiting, context length violations, and invalid request parameters.
    \item \texttt{PassportOAuthErrorHandler} — handles Passport OAuth authentication errors such as invalid or expired tokens, profile retrieval failures, and authorization denial by the user.
    \item \texttt{RateLimitErrorHandler} — handles application rate limiting errors by detecting excessive client requests and returning an appropriate rate limit response.
  \end{itemize}

\paragraph{The Registry}
The \texttt{errorHandlerRegistry} is an ordered array of handler instances. The middleware iterates through it and delegates to the first handler whose \texttt{canHandle()} returns \texttt{true}. Order is 
significant: specific handlers must precede general ones, and \texttt{UnexpectedErrorHandler} must always be last.

\subsubsection{The Error Handler Middleware}
The middleware contains no error-type-specific logic. It iterates the registry and delegates entirely:

\begin{lstlisting}
export const errorHandler = (err, req, res, next) => {
  const handler = errorHandlerRegistry.find(h => h.canHandle(err));
  if (handler) return handler.handle(err, res);
};
\end{lstlisting}

\subsubsection{Extending the System}
Adding support for a new third-party library error requires two steps:

\begin{enumerate}
    \item Implement the \texttt{IErrorHandler} interface in a new handler class, defining \texttt{canHandle()} for detection and \texttt{handle()} for the response. \item Register the new handler in \texttt{errorHandlerRegistry} before \texttt{AppErrorHandler}.
\end{enumerate}

For known, expected application failures such as not found or unauthorized, no new handler is needed. Throwing \texttt{new AppError(message, statusCode)} anywhere in the application is sufficient, as \texttt{AppErrorHandler} handles it automatically.

\subsubsection{Design Evaluation}
The error handling mechanism offers clear separation of concerns — each handler owns both the detection and response logic for exactly one error category. The middleware is fully decoupled from third-party libraries, 
depending only on the \texttt{IErrorHandler} abstraction. Each handler is independently testable, with unit tests verifying detection and response logic in isolation. Adding a new error type extends the system 
without modifying any existing code.
